
\subsubsection{6.1 Generation-Free Detection}

For commission-dominated tasks, post-hoc NLI scoring on logged (context, response) pairs yields AUROC = 0.644--0.709 at inference time cost effectively zero beyond the NLI classifier forward pass. The generation-free margin over SelfCheckGPT-NLI is +0.229 on dialogue and +0.114 on QA under this evaluation. The magnitude of the SelfCheckGPT disadvantage is partially attributable to evaluation on base LLaMA-3-8B rather than instruction-tuned variants; the true margin under fair comparison is not established by this study.

\subsubsection{6.2 The Commission/Omission Boundary}

The task-level pattern in Table 1 --- AUROC 0.709/0.644 for commission tasks, 0.530 for the omission task --- is consistent with a structural account: NLI models trained to detect textual contradiction are suited to detecting commission-type hallucinations and unsuited to detecting omission-type. Prior work (SummaC, TRUE) observed task-level variation in NLI-based consistency checking without providing a unified mechanistic account. This paper presents, to our knowledge, the first AUROC-based quantitative characterization of this boundary across three tasks. The commission/omission distinction predicts where NLI-based generation-free detection is viable without task-specific tuning. Hybrid approaches --- NLI contradiction scoring for commission detection combined with coverage-based metrics (BERTScore recall, ROUGE recall) for omission detection --- are a natural direction.

\subsubsection{6.3 Limitations}

\textbf{L1 (Summarization scope, High)}: AUROC = 0.530 on HaluEval-Summarization, near chance and near the structural ceiling (0.52). The mismatch between NLI contradiction detection and omission-type hallucinations is architectural; it is not correctable by threshold adjustment, model scaling, or minor configuration changes.

\textbf{L2 (QA near-miss, Medium)}: AUROC = 0.644 on QA falls 0.006 below the pre-specified 0.65 target. Cohen's d = 0.779 indicates strong practical discrimination, but the result does not meet the original hypothesis threshold.

\textbf{L3 (Incomplete ablations, High)}: Design-choice ablation experiments (h-m2 through h-m4) evaluating net-contradiction framing advantage, sentence-level aggregation benefit, and dialogue windowing were not executed. The reported AUROC values characterize the complete configuration but cannot be attributed to individual components. The current results confirm the configuration functions but cannot distinguish whether simpler configurations would achieve comparable performance.

\textbf{L4 (Single model, Medium)}: All results derive from one frozen NLI model (\texttt{cross-encoder/nli-deberta-v3-large}). Generalizability to other NLI architectures (roberta-large-mnli, bart-large-mnli) is not established.

\textbf{L5 (Benchmark ecological validity, Medium)}: HaluEval labels are generated by GPT-3.5 perturbations rather than naturalistic LLM outputs. Detection performance on real-world hallucination distributions may differ.

\textbf{L6 (KL gate criterion, Low)}: The pre-specified KL > 0.05 threshold is not calibrated to task-level differences in score distribution variance. QA's KL = 0.035 reflects distribution compression in short-context tasks, not absence of mechanism signal. Task-specific thresholds or effect-size-based criteria are more appropriate for multi-task evaluation.

\subsubsection{6.4 Broader Impact}

Post-hoc NLI detection is applicable to any logged (context, response) pair collection without additional LLM inference. For dialogue and QA applications where hallucinations are predominantly commission-type, this provides a low-cost monitoring path. The commission/omission framework also provides guidance for cases where this approach is not appropriate: high-stakes summarization applications require omission-aware detection that contradiction scoring cannot provide.

The h-m1 gate failure and h-m1 SELF\_MODIFY outcome are disclosed for transparency: the original mechanism gate criterion (KL > 0.05 on all tasks) was not met by QA (KL = 0.035), and Wilcoxon rank-sum and Cohen's d were subsequently used as primary mechanism indicators. This should be noted when interpreting the mechanistic claims.
